\documentclass{article}
\usepackage[UTF8]{ctex}
\usepackage{amsmath,amssymb,amsfonts}
\usepackage{enumitem}
\usepackage{geometry}
\geometry{a4paper, margin=1in}
% === 错题本排版优化（不参与题目标识） ===
\setlength{\parindent}{0pt}
\setlength{\parskip}{0.35em plus 0.10em minus 0.08em}
\setlength{\emergencystretch}{3em}
\setlist{itemsep=0.25em, topsep=0.40em, parsep=0pt}
\allowdisplaybreaks[2]
% === 错题本排版优化结束 ===
\begin{document}
\section*{华南理工大学2025年高等代数真题}
\section*{146．华南理工大学2025年高等代数试卷}

\subsection{题目1}
1．（15 分）设 \(f(x)\) 是非零整系数多项式，证明：

\begin{itemize}
  \item[（1）] 对任意整数 \(a\) ，存在整系数多项式 \(q(x)\) 和整数 \(r\) ，使得 \(f(x)=(x-a) q(x)+r\) 。
  \item[（2）] 如果有两两不同的整数 \(a, b, c\) ，使得 \(|f(a)|=|f(b)|=|f(c)|=1\) ，则 \(f(x)\) 没有整数根．
\end{itemize}

\subsection{题目2}
2．（20 分）设矩阵 \(A=\left(\begin{array}{cccc}a_{11} & a_{12} & \cdots & a_{1 n} \\ a_{21} & a_{22} & \cdots & a_{2 n} \\ \vdots & \vdots & & \vdots \\ a_{n 1} & a_{n 2} & \cdots & a_{n n}\end{array}\right)\) ，且

\[
a_{i 1}+a_{i 2}+\cdots+a_{i n}=0, a_{1 j}+a_{2 j}+\cdots+a_{n j}=0(i, j=1,2, \cdots, n) .
\]

证明：\(A\) 的所有元素的代数余子式都相等。

\subsection{题目3}
3．（20 分）\(\lambda\) 取何值时，线性方程组

\[
\left\{\begin{array}{l}
-\lambda x_{1}+2 x_{2}-2 x_{3}=1 \\
2 x_{1}+(3-\lambda) x_{2}-4 x_{3}=2 \\
-2 x_{1}-4 x_{2}+(3-\lambda) x_{3}=-\lambda-3
\end{array}\right.
\]

有唯一解，无解，有无穷多解？在有无穷多解时求通解（用基础解系表出）。

\subsection{题目4}
4．（20 分）设 \(N\) 是 \(n\) 阶幂零矩阵，即存在自然数 \(k\) ，使得 \(N^{k}=O\) 。

\begin{itemize}
  \item[（1）] 证明：\(N+E\) 可逆，其中 \(E\) 是 \(n\) 阶单位矩阵。
  \item[（2）] \(A\) 是 \(n\) 阶可逆矩阵，且 \(A N=N A\) ，证明：\(A+N\) 可逆．
\end{itemize}

\subsection{题目5}
5．（20 分）设 \(A\) 是实数域 \(\mathbb{R}\) 上的 \(n\) 阶方阵，\(\alpha \in \mathbb{R}^{n}\) 是 \(n\) 维列向量．证明：\(A-\alpha \alpha^{\mathrm{T}}\) 正定当且仅当 \(A\)正定。且 \(\alpha^{\mathrm{T}} A^{-1} \alpha<1\) ．

\subsection{题目6}
6．（15 分）设 4 阶方阵 \(A=\left(\begin{array}{cccc}2 & -1 & 1 & 0 \\ -1 & 2 & -1 & 0 \\ 1 & -1 & 2 & 0 \\ 2 & a & b & 1\end{array}\right)\) 在复数域上可对角化．

\begin{itemize}
  \item[（1）] 求 \(A\) 的特征多项式．
  \item[（2）] 确定 \(a, b\) 的值。
  \item[（3）] 求可逆矩阵 \(P\) ，使得 \(P^{-1} A P\) 为对角阵。
\end{itemize}

\subsection{题目7}
7．（20 分）设

\[
J=\left(\begin{array}{llll}
\lambda & 1 & & \\
& \lambda & \ddots & \\
& & \ddots & 1 \\
& & & \lambda
\end{array}\right)
\]

是数域 \(P\) 上的 \(n\) 阶若尔当块矩阵，\(V=\left\{A \in P^{n \times n} \mid A J=J A\right\}\) 。证明：

\begin{itemize}
  \item[（1）] \(V\) 是线性空间．
  \item[（2）] \(\operatorname{dim} V=n\) ．
\end{itemize}

\subsection{题目8}
8．（20 分）设 \(V\) 是 \(n\) 维线性空间，\(V_{1}, V_{2}\) 是 \(V\) 的子空间，且 \(\operatorname{dim} V_{1}+\operatorname{dim} V_{2}=n\) ．证明：存在 \(V\)上的线性变换 \(\mathcal{A}\) ，使得 \(\mathcal{A} V=V_{1}, \mathcal{A}^{-1}(0)=V_{2}\) ．
\end{document}
